%% 第一章

\chapter{绪论}
\chaptermark{绪论}

\section{课题研究背景与意义}

近年来，复杂环境下的自主机器人技术快速发展，移动机器人与协作机械臂等智能体常需在高动态、非结构化等复杂工况下执行任务。面对未知且高风险的真实物理环境，依靠真实世界进行智能体感知、决策与控制算法的训练不仅成本极为高昂，且存在巨大的系统损毁风险。因此，构建高保真的虚拟仿真场景进行算法的验证与部署前评估，已成为一项关键途径。然而，传统的物理仿真环境通常依赖繁琐的人工建模，在复杂场景生成效率、物理参数映射的细粒度以及场景的视觉-物理时空一致性上存在显著不足，导致虚拟环境与真实环境差异较大，难以满足泛化智能体高效、高精度的 Sim2Real（仿真到现实）迁移需求。

随着以神经辐射场（NeRF）\cite{nerf}和三维高斯溅射（3DGS）\cite{3dgs}为代表的三维表征技术与视觉大模型的突破，基于单图、多视角或文本序列的高效场景外观生成技术取得了长足进步。但这些方法多侧重于静态几何与光学分布，缺乏内在的质量、弹性、密度等力学本构属性，无法在虚拟场景中产生符合真实物理规律的动态交互。部分前沿研究尝试将连续介质力学模拟（如物质点法 MPM）融合进基于散点（如 3DGS）的表达中\cite{physgaussian}，但大多局限于单一材质或无复杂背景干扰的简单物体。当面对包含复杂地貌、探索车支撑结构等多材质耦合的庞大环境时，由于缺乏场景级的部件辨别与材质自动识别机制，往往面临材质分割模糊、物理先验缺失等问题\cite{gaussianproperty}；更为严重的是，大规模静态背景在仿真中牵连整个网格系统，易导致网格刚体化和梯度计算发散等瓶颈。

针对现有复杂场景仿真中物理真实性不足与生成效率低下等痛点，本文以弥合视觉场景生成与复杂物理动力学演化之间的鸿沟为核心目标，探索并构建了物理知识嵌入的复杂场景生成与仿真系统（PhysLucid）。本课题旨在通过集成视觉分割大模型（SAM）\cite{sam}、视觉语言模型（VLM）与三维特征溅射（Feature-Splatting）\cite{featuresplatting}，在三维空间中自动化地完成材质属性分布与本构模型的识别及蒸馏，从而构建出适用于复杂非结构化环境的高保真物理交互虚拟场景。

本课题的研究兼具理论价值与工程应用意义。在理论与算法层面，系统创新性地打通了"多模态语义理解——三维特征蒸馏——多材质物理扩展演化"的技术链路；通过分离庞大静态地貌与动态前景实体的解耦仿真架构，打破了经验手动调参的壁垒，有效解决了大规模背景牵制导致的物理交互失效瓶颈。在工程与应用层面，课题构建的生成式仿真系统为复杂场景下部署自主智能体搭建了高置信级别的闭环仿真验证平台。此项研究以可计算、可微调的弹塑性力学机制赋能传统的三维静态渲染，旨在大幅缩减智能体算法在 Sim2Real 迁移过程中的性能损耗。其成果将显著提升协作智能体在复杂场景中的任务执行成功率与系统鲁棒性，为具身智能领域的高精度物理仿真研究提供有效的技术支撑。

\section{国内外研究现状}

针对"物理知识嵌入的复杂场景生成与仿真"这一交叉前沿领域，学术界和工业界在三维重建、多模态语义理解以及连续体力学仿真等层面取得了丰硕的成果。结合本课题的探索方向，现对相关核心技术进展梳理如下：

\subsection{基于三维高斯的场景重建}
自神经辐射场（NeRF）\cite{nerf}提出以来，基于隐式神经表示的三维新视角合成获得了突破性进展。在此基础上，由 Kerbl 等人开创的三维高斯溅射（3DGS）\cite{3dgs}技术迅速崛起，其将场景显式表征为数以百万计的各向异性三维高斯椭球散点，利用可微光栅化技术在保证极高保真度的同时，全面突破了此前隐式表征的渲染速率瓶颈，实现了高分辨率的实时渲染。在此基础上，Deformable 3D Gaussians\cite{deformable3dgs} 进一步将静态高斯扩展至时序维度，通过在规范空间中学习形变场，使得高斯原语能够随帧间运动发生持续性形变，为后续将 3DGS 与物理仿真衔接提供了动态表征的理论基础。然而，从实际任务需求出发，无论是 NeRF 还是传统的 3DGS，其本质都在于构建高度一致的视觉表象——得到的虚拟实体大多仅具备光影与几何包络，缺乏质量、弹性本构、接触体积等基础力学属性，无法直接支持如机械臂触碰或移动机器人行驶等真实的物理阻力交互。

\subsection{三维场景生成}
在三维重建技术日趋成熟的基础上，从文本或单图直接生成完整三维场景成为近年来的研究热点。DreamFusion\cite{dreamfusion} 率先提出以 2D 扩散模型作为先验、通过分数蒸馏采样（SDS）将文本描述转化为三维表征的范式，开创了文本到三维的生成先河。在此基础上，LucidDreamer\cite{luciddreamer} 将生成对象从单个物体扩展至完整场景，利用 3DGS 显式表征实现了无域限制的高质量场景生成。Director3D\cite{director3d} 进一步引入了相机轨迹扩散与场景几何的联合建模，显著提升了多视角一致性与空间结构的合理性。FlashWorld\cite{flashworld} 则在生成效率上取得突破，将高质量三维场景的构建时间压缩至秒级。WonderWorld\cite{wonderworld} 探索了单图驱动的交互式三维场景生成路径，允许用户从单一参考图像出发逐步扩展和编辑虚拟环境。然而，当前三维场景生成方法大多聚焦于视觉外观的逼真度，极少考虑场景内各部件所应具备的物理本构属性，导致生成结果无法直接支撑后续的力学仿真与动态交互。

\subsection{语义蒸馏与特征感知体系}
传统物理仿真环境的参数配置极其依赖人工设定与经验调优。为了赋予三维场景自动化的级联属性识别能力，当下的前沿范式开始借力视觉分割基础大模型（如 SAM\cite{sam}）及视觉-语言多模态大模型（VLM）。LERF\cite{lerf} 率先将 CLIP 语义特征嵌入神经辐射场，实现了三维空间内的零样本语言查询，开创了视觉语言模型与三维表征融合的先河；Feature-Splatting\cite{featuresplatting} 则在此基础上将特征蒸馏引入 3DGS 显式表征，通过可微光栅化实现语言驱动的物理场景编辑与合成。在三维分割层面，SegAnyGaussian\cite{seganygaussian}（SAGA）通过在 3D 高斯上附加尺度门控亲和力特征，利用 SAM 的二维掩码知识蒸馏，实现了毫秒级的三维可提示分割；Gaga\cite{gaga} 则构建了 3D 感知记忆库机制，在无需连续视角假设的前提下实现跨视角的高斯分组与场景理解。GaussianProperty\cite{gaussianproperty} 进一步探索了利用多模态大模型为 3D 高斯推断并赋予物理属性（如杨氏模量）的可行路径，为物理仿真提供了自动化材质标注的初步方案。但以往的特征提取研究多数停留在视觉效果编辑与增强检索层面，鲜有研究将大模型赋予的自然属性与语义完整跨越至底层的弹性模量等物理计算域。这种"文本理解——视觉部件识别——物理常识锚定"链条之间，仍存在亟待填补的断层。

\subsection{物理知识嵌入三维场景的技术路线}
物理知识嵌入（Physics-Knowledge Embedding）指将质量、弹性、摩擦、碰撞响应等物理属性赋予纯视觉三维表征，使其不仅"看起来真"，更能"动起来真"。围绕这一目标，学术界近年来涌现出多条各具特色的技术路线，在物理保真度、计算效率与自动化程度之间形成了丰富的频谱分布。

\textbf{连续介质力学与三维表征的耦合。}该路线将物质点法（MPM）、有限元法（FEM）等传统数值求解器与 NeRF、3DGS 等三维视觉表征直接融合，是目前物理保真度最高的范式。PhyRecon\cite{phyrecon} 首次将可微渲染与 FEM 耦合至统一隐式表面框架，在重建过程中联合优化外观与力学稳定性。PhysGaussian\cite{physgaussian} 将 3DGS 显式散点直接映射为 MPM 质点，建立了高斯溅射到连续介质力学的无缝通路。PhysDreamer\cite{physdreamer} 进一步将可微 MPM 管线推广至多材质交互场景，实现了物体间碰撞、挤压等物理响应的无监督学习。在参数优化层面，DreamPhysics\cite{dreamphysics} 以 SVD 视频扩散模型作为教师信号，通过 SDS 损失将运动规律蒸馏为物理参数；PhysGen3D\cite{physgen3d} 探索了单图输入下微型可交互场景的生成路径。该类方法的优势在于物理严格性——遵循守恒定律与本构方程，仿真结果具有物理可解释性；其挑战在于求解器计算代价高昂，且在复杂多材质、大规模背景场景中面临网格刚体化与梯度发散等稳定性瓶颈。

\textbf{可微刚体动力学与多体仿真。}与连续介质力学并行的另一条路线是刚体动力学框架与视觉表征的耦合。新一代可微刚体物理引擎将接触力学与碰撞响应完全 GPU 化并暴露梯度接口，使梯度能够穿过物理仿真步骤从渲染像素反传至场景参数，为机器人抓取、物体堆叠等刚体主导的场景提供了兼具视觉真实性与物理一致性的仿真方案。然而，当前可微刚体路线在应对大变形、断裂等连续介质力学问题时尚属空白，与本文聚焦的弹性体、织物等柔性材质仿真呈互补关系。

\textbf{数据驱动的神经物理仿真器。}随着图神经网络（GNN）与 Transformer 等深度学习架构的成熟，学术界开辟了完全绕过传统数值求解器的"端到端神经物理仿真"路线。该类方法将物理系统建模为粒子/网格节点图，通过消息传递或自注意力机制隐式学习节点间作用力，一步到位地从当前状态预测下一时刻的动力学演化，在流体动力学和布料仿真等任务上展现出超越传统求解器的速度优势。然而，该类"黑箱"方法的泛化能力严重依赖训练数据分布，遇到训练分布外的材质或边界条件时往往产生非物理的预测，且难以与基于明确物理参数（如杨氏模量 $E$、泊松比 $\nu$）的传统工程设计流程对接——后者正是可微传统求解器路线的核心优势，也是本文选择在 MPM 显式物理框架内进行参数寻优的深层原因。

\textbf{视频生成模型作为隐式物理先验。}与上述显式赋予物理规则的方法不同，一条极具潜力的新兴路线是利用大规模视频生成模型内部蕴含的"物理常识"作为隐式监督信号。这类方法的核心洞察在于：在数亿级视频数据上训练的扩散模型不仅学会了物体的外观分布，更隐式习得了符合现实世界统计规律的运动模式。DreamPhysics\cite{dreamphysics} 率先将这一思想系统化，以 SVD 作为"物理评分器"通过 SDS 梯度优化 MPM 参数。PhysDreamer\cite{physdreamer} 进一步探索了利用视频扩散先验统一驱动静态 3D 资产产生符合物理常识的动态交互。相较于传统依赖真实传感器数据或人工标注的物理标定流程，视频扩散先验路线的最大优势在于无需任何真值物理参数标注即可实现自监督优化；但其挑战在于扩散模型的物理合理性受限于训练数据的统计偏差，且 SDS 梯度的高方差特性导致优化过程收敛缓慢。本文工作正是在上述路线交叉点上展开——以连续介质力学（MPM）为物理底座，以视频扩散先验（SVD + SDS）为优化驱动力，以多模态语义蒸馏（VLM + Feature-Splatting）为材质先验注入机制，形成一条融合"显式数值求解器 + 隐式扩散先验 + 多模态语义引导"的混合物理知识嵌入路径。

\section{论文主要工作内容}

针对目前三维场景仿真物理交互性不足、复杂背景导致网格刚体化发散以及多材质环境物理赋值严重依赖人工等痛点，本文设计并实现了一套物理知识嵌入的复杂场景生成与仿真系统（PhysLucid）。该系统旨在贯通三维视觉快速重建与力学规律高保真演化的技术链路，主要研究内容包含以下四个核心部分：

第一，面向物理仿真的三维场景构建与材质初始化。以文本或参考图像作为输入，引入 FlashWorld\cite{flashworld} 零样本生成多视角图像与相机位姿，快速构建三维高斯场景雏形。针对复杂场景中的遮挡混叠问题，采用 Depth-Anything-V2\cite{depthanythingv2} 估计逐像素深度并实施深度分层（Depth Split），将前景物体与背景在深度维度上解耦。在此基础上，利用 SAM\cite{sam} 对分层后的各深度图层进行自动分割，获取独立的物体掩码区域；进而调用视觉语言模型（VLM）对每个孤立区域进行材质推理与分类，最终将材质语义映射为杨氏模量、泊松比与密度的初始物理属性范围，为后续梯度寻优提供约束先验。

第二，基于多源语义的三维特征溅射蒸馏（Feature-Splatting）\cite{featuresplatting}。该模块承担两个层次的蒸馏任务：其一，提取二维图像空间中的 CLIP\cite{clip} 与 DINO\cite{dino} 视觉语义特征，通过可微 $\alpha$-blending 光栅化将对象级/零件级视觉特征联合蒸馏至每个三维高斯原语；其二，将第三章 VLM 推理得到的材质语义标签通过独立语义解码器同步蒸馏至高斯空间，并引入逐高斯语义一致性损失、独立语义特征库与时延损失调度等机制，确保视觉语义与物理材质分类在三维空间中的双重对齐。

第三，构建耦合前景动态剥离机制的解耦力学演化系统。面对大尺度地形地貌等多构件引发的系统算力崩溃问题，研究自然语言驱动的特征分割解绑方案。引入前景分割（Mark-Foreground）背景软冻结逻辑，基于文本向量与自适应预设多策略分割框架精准屏蔽无关大体积背景干扰，并辅以 GPU 加速连通分量分析（CCA）与 SAM 掩码反投影作为补充链路；利用物质点法（MPM）\cite{physgaussian}对筛选出的前景高斯散点进行力学降维求导，实现场景内多模态弹性体平滑逼真的物理再现。

第四，基于视频扩散先验蒸馏的物理参数反演（DreamPhysics）。本系统摒弃了复杂且繁琐的手工参数标定，在末端集成 DreamPhysics\cite{dreamphysics} 自适应物理寻优模块：该模块以 Stable Video Diffusion（SVD）图像到视频扩散模型作为隐式物理先验——SVD 在数亿级视频数据上习得的运动规律被用作"伪真值"教师信号，通过分数蒸馏采样（SDS）将 MPM 仿真渲染帧与 SVD 预期运动之间的分布偏差转化为弹性模量 $E$ 与密度 $\rho$ 的优化梯度，经 Warp 全可微物理管线反传至每一物质点。在寻优过程中，系统集成了运行时 CFL 自适应时间步长、逐粒子杨氏模量钳位、体素材质正则化与按材质类型梯度归一化等多重数值稳定性保障机制。该模块将 SVD 视频扩散模型内含的物理常识蒸馏为可解释的连续介质力学参数，实现了从"看视频懂物理"到"算物理像物理"的跨模态知识迁移。系统最终搭建多种对比与消融实验，定性与定量验证了其在改善仿真物理真实性方面的有效性与鲁棒性。

\section{论文组织结构}

本文的组织结构安排如下：

第一章：绪论。介绍了复杂场景物理仿真在高动态工况下的研究背景与紧迫需求，系统梳理了三维表征与重建、场景生成、多模态语义蒸馏、以及物理知识嵌入三维场景的四条主流技术路线（连续介质力学耦合、可微刚体动力学、数据驱动神经仿真器、视频生成隐式物理先验）的国内外研究现状，并在此基础上明确了本文融合"显式求解器 + 隐式扩散先验 + 多模态语义引导"混合路径的核心研究内容。

第二章：相关技术基础。系统构建了本文框架依托的理论基石，沿着"三维场景表征—视觉语义理解—力学仿真与视频生成"三条主线，分别阐述神经辐射场与三维高斯溅射的渲染理论、SAM/CLIP/DINO/VLM 等视觉基础模型的架构原理、以及物质点法（MPM）数值求解框架与 Stable Video Diffusion（SVD）扩散模型的生成机理。

第三章：面向物理仿真的三维场景构建与材质初始化。详解从文本或参考图像到物理可仿真三维场景的全链路预处理架构，包括 FlashWorld 零样本多视角生成与相机轨迹后处理、基于 Depth-Anything-V2 的三层深度重切分（Depth Split）遮挡消解策略、SAM 自动网格提示分割与掩码筛选、Qwen3-VL 多模态材质识别与分类、以及基于复合 ID 编码（composite\_id = label×10 + grade）的材质等级体系与物理常数经验映射。

第四章：基于三维高斯的语义特征蒸馏与前景背景分离。阐述如何通过扩展高斯原语挂载视觉语义与材质语义两组独立高维向量，并引入 YOLO-MobileSAMv2 检测引导的特征预提取管线；详述 CLIP 与 DINO 双流特征蒸馏策略、DINO 递增预热机制、独立语义特征库、逐高斯语义一致性损失、特征引导的语义平滑约束与时延损失调度策略，以及材料高斯点密度上限控制等训练优化与显存安全保障。

第五章：基于视频扩散先验的物理参数反演与连续体仿真。论述从语义标签到 MPM 求解器类型的物理参数自动映射（含求解器类型匹配、少数类中值钳位、CFL 硬 E 上限与自适应网格尺寸），详述五阶段多源证据融合的前景/背景分离管线（Mark-Foreground），以及 DreamPhysics 模块中 SVD 视频扩散先验引导的物理参数 SDS 梯度反演过程，并分析运行时 CFL 自适应、逐粒子钳位、体素材质正则化与按材质类型梯度归一化等数值稳定性保障。

第六章：平台实验结果与系统分析。从系统级验证角度，构建包含帽子、花瓶、雕塑等多类复杂场景的评估数据集，对深度分层、材质分类、前景分割、物理参数寻优等关键模块进行定量消融实验与定性可视化分析，评估多材质耦合动态响应与极端工况下的系统鲁棒性。

第七章：总结与展望。对全文研究工作进行系统化总结，反思当前系统在材质识别粒度、仿真稳定性边界与计算效率等方面的局限，并在此基础上对复杂场景高保真物理仿真的未来研究方向提出展望。
